\documentclass[a4paper,11pt]{article}
\usepackage[utf8]{inputenc}
\usepackage[T1]{fontenc}
\usepackage[french]{babel}
\usepackage[left=2cm,right=2cm,top=2cm,bottom=2cm]{geometry}
\usepackage{xcolor}
\usepackage{hyperref}
\usepackage{fontawesome5}
\usepackage{parskip}

% --- Couleurs CEA ---
\definecolor{primary}{RGB}{0, 102, 153}
\hypersetup{colorlinks=true, urlcolor=primary}

\begin{document}

% --- EXPÉDITEUR ---
\noindent
\begin{minipage}[t]{0.5\textwidth}
    \textbf{Loïc Aron MBASSI EWOLO}\\
    Élève-Ingénieur Bac+5 -- Double diplôme ENSEA\\[3pt]
    \faMapMarker\ Cergy, France\\
    \faPhone\ (+33) 7 59 52 33 98\\
    \faEnvelope\ \href{mailto:loicaron.mbassiewolo@ensea.fr}{loicaron.mbassiewolo@ensea.fr}
\end{minipage}
\hfill
% --- DATE ---
\begin{minipage}[t]{0.4\textwidth}
    \raggedleft
    Cergy, le 2 mars 2026
\end{minipage}

\vspace{1.2cm}

% --- DESTINATAIRE ---
\hfill
\begin{minipage}[t]{0.5\textwidth}
    \raggedleft
    \textbf{CEA Gramat}\\
    Service Recrutement\\
    Gramat, France
\end{minipage}

\vspace{1cm}

\noindent\textbf{Objet :} Candidature au stage -- Développement d'un outil d'automatisation : validation expériences-calculs

\vspace{0.8cm}

Madame, Monsieur,

Développer un outil Python d'automatisation pour confronter résultats de simulation et données expérimentales multi-physiques : ce sujet conjugue deux compétences que je mets en pratique depuis plusieurs années, le développement d'outils logiciels et l'analyse de données scientifiques. En double diplôme à l'ENSEA après un cursus d'ingénieur en Mathématiques Appliquées à l'École Polytechnique de Yaoundé, je souhaite mettre cette expertise au service de l'équipe de dynamique rapide du CEA Gramat.

Mon expérience la plus directement liée à ce stage est la conception d'un outil open source (Laravel API Generator, +30 étoiles GitHub) qui automatise la génération de code backend depuis des diagrammes UML/XML. Ce projet m'a confronté aux mêmes enjeux que votre sujet : définir un cahier des charges en concertation avec les utilisateurs, parser et exploiter une base de données structurée, et produire une documentation technique claire pour garantir la maintenabilité. L'outil est aujourd'hui utilisé par d'autres développeurs grâce à une documentation rigoureuse et un déploiement sur Packagist.

Sur le volet scientifique, mon stage au MILA (Institut Québécois d'IA) m'a formé à l'analyse mathématique de modèles (étude du phénomène de Grokking sur des MLP, reproduction d'expériences SOTA en Python, rédaction de rapports de synthèse). Le projet UROP sur l'analyse d'électrocardiogrammes, mené en collaboration avec des mentors Google DeepMind, m'a appris à confronter les résultats d'un modèle avec des données expérimentales de référence en utilisant des techniques de traitement du signal et d'analyse de séries temporelles. Ces expériences m'ont donné le réflexe de valider systématiquement les résultats numériques par rapport aux observations expérimentales.

Je suis disponible dès avril 2026 pour une durée de 4 à 6 mois et mobile pour rejoindre le site de Gramat. La perspective d'une suite après le stage renforce ma motivation à livrer un outil robuste et bien documenté. Je reste à votre disposition pour un entretien afin de discuter de ma candidature.

Je vous prie d'agréer, Madame, Monsieur, l'expression de mes salutations distinguées.

\vspace{0.8cm}

\textbf{Loïc Aron MBASSI EWOLO}\\
\textit{Élève-Ingénieur ENSEA / Polytechnique Yaoundé}

\end{document}
